\section{A2 EEG 语义融合模型}

本章详细介绍本课程设计中的关键定量模型——A2 temporal-spectral-spatial EEG 语义融合模型。该模型的基本想法是把 EEG 信号当作视觉语义的增强来源：当输入图像发生退化时，利用同一时刻采集的 EEG 数据，在 CLIP 语义空间中对退化图像的表示进行补偿和校准，以提高语义分类的准确率。该模型全称为 Temporal-Spectral-Spatial Transformer，简写为 A2，在 CLIP 语义空间内同步提取 EEG 的时间维度、频谱维度和空间维度特征，并通过门控残差融合机制来修正退化图像的语义表示。

\subsection{设计动机}

EEG 本质上是一个高维时间序列。在本课程设计所用的数据集中，每一条 EEG trial 的形状是 $[64, 250]$，即 64 个电极通道在 250 个时间采样点记录的电压幅值。该数据蕴含三个维度的信息。时间维度刻画了刺激呈现后脑电活动随时间的演变，包括早期视觉诱发响应（P100/N170）和晚期的语义相关成分（P300）。频谱维度对应不同频段（$\alpha$、$\beta$、$\gamma$）的振荡能量分布，诸多研究表明这些频段分别与视觉加工、注意分配以及语义理解等过程存在关联。空间维度反映了 64 个电极在头皮上的几何排布，隐含了不同脑区之间的功能连接关系。

大多数单一的编码器架构只能较好地利用上述三个维度中的一到两个。例如，时序卷积网络擅长捕捉局部时间模式，却难以建模通道间的空间关系；基于频谱的 CNN 可以从时频域提取特征，但会丢失原始时间拓扑，且跨频段交互建模能力有限；Transformer 直接作用于原始时序虽能捕获长距离依赖，但缺少显式的频谱建模。A2 的策略是为每个维度设计一个专用的编码分支，各分支采用适配该维度特性的结构，再通过一个后融合层完成跨维度的信息整合。在模型容量相对有限的前提下，这种显式分离再融合的设计相比单纯堆叠同一类编码器更有效率。

在具体结构选型上，A2 吸收了 EEG 信号处理领域的若干已有成果。ConvTransformer 在 raw EEG 时序建模中已被验证有效——其卷积 stem 提供局部平滑，Transformer 层则负责全局依赖关系建模。基于 STFT 的频谱 CNN 在人脑 EEG 情感识别和运动想象分类任务中已有成功先例。通道空间注意力机制也在基于 EEG 的脑机接口中被用于学习电极之间的功能连接。A2 的贡献在于将这三个方向统一至一个端到端可训练的框架内，并针对视觉语义解码这一具体任务进行了联合优化。

\subsection{A2 EEG Encoder 结构}

A2 编码器由三个并行的分支模块和一个后融合 Transformer 构成。其整体结构如图~\ref{fig:a2-encoder} 所示。

\placeholderfigure{A2 temporal-spectral-spatial EEG encoder 结构图}{a2-encoder-placeholder}{5.0cm}

\textbf{Temporal 分支}：该分支采用多尺度 ConvTransformer 设计。四个 Conv1d 子分支并行运行，卷积核尺寸分别为 $k \in \{3, 7, 15, 31\}$，从不同尺度提取时间特征。各分支的输出拼接后经 $1\times1$ 卷积进行通道混合，再由 4 层 Transformer Encoder 建模长程时序依赖。最后通过 learnable attention pooling 和 MLP 投影得到 512 维的 temporal embedding。这种多尺度策略允许模型同时捕获毫秒量级的瞬态变化和持续数百毫秒的慢变趋势。

\textbf{Spectral 分支}：该分支基于可微分的短时傅里叶变换构建频谱 CNN。对 EEG 信号施加 STFT（窗长 64、hop length 16），将幅度谱取对数后逐通道标准化，再用两层 2D CNN 提取时频域中的二维模式。经全局平均池化和 MLP 投影后输出 512 维的 spectral embedding。

\textbf{Spatial 分支}：该分支先将 EEG 信号转置，使每个通道的时间序列经由线性层映射为 hidden\_dim 维的 token，再利用 learnable attention pooling 进行加权聚合，最后通过 MLP 输出 512 维的 spatial embedding。该分支无需预定义通道图结构，依靠自注意力机制自动学习各电极之间的功能连接关系。

\textbf{后融合}：三个分支各自输出 512 维向量，拼接为 1536 维后送入 2 层 Transformer Encoder 进行跨分支交互融合。融合 Transformer 使用可学习的 branch type embedding 来标注各特征的来源分支。最终经 attention pooling 和 MLP 投影器输出 512 维的融合 EEG embedding $f_{\mathrm{eeg}}$。

\subsection{CLIP 文本原型与语义融合}

CLIP 文本原型的构造方式与第 2 章一致：为全部 40 个类别分别构建 prompt ``a photo of a \{class\_name\}''，经冻结的 CLIP text encoder 编码后得到归一化的文本原型向量。在语义预测阶段，A2 输出的融合 EEG embedding $f_{\mathrm{eeg}}$ 与退化图像的视觉 embedding $f_{\mathrm{img}}$ 通过 gated residual fusion 进行混合：

\begin{equation}
  f_{\mathrm{fused}} = f_{\mathrm{img}} + \sigma\left(g\left([f_{\mathrm{img}}; f_{\mathrm{eeg}}]\right)\right) \odot \Delta(f_{\mathrm{eeg}})
\end{equation}

其中 $g(\cdot)$ 是一个 2 层 MLP 构成的门控网络，$\sigma$ 为 sigmoid 激活函数，$\Delta(\cdot)$ 是对 EEG 做残差变换的网络。门控机制根据图像与 EEG 的联合信息动态调控 EEG 的注入强度：当图像质量较好时门控系数趋近于 0，从而削弱 EEG 的干预；当图像退化严重时系数增大，使 EEG 信息得到更充分的利用。

融合向量 $f_{\mathrm{fused}}$ 随后与 40 个文本原型计算余弦相似度，经 softmax 得到各类别的预测概率。图~\ref{fig:a2-fusion} 展示了完整的流程。

\placeholderfigure{A2 semantic fusion 流程图}{a2-fusion-placeholder}{4.0cm}

\subsection{训练设计}

A2 语义融合模型以交叉熵分类损失作为主训练目标。对于融合嵌入 $f_{\mathrm{fused}}$ 与 40 个文本原型 $\{p_c\}_{c=1}^{40}$，将余弦相似度经 softmax 转化为类别概率：

\begin{equation}
  \mathcal{L}_{\mathrm{CE}} = -\frac{1}{B} \sum_{i=1}^{B} \log \frac{\exp(\tau^{-1} \cos(f_{\mathrm{fused},i}, p_{y_i}))}{\sum_{c=1}^{K} \exp(\tau^{-1} \cos(f_{\mathrm{fused},i}, p_c))}
\end{equation}

式中 $\tau$ 是可学习的温度参数，$y_i$ 表示第 $i$ 个样本的真实类别。

参数更新仅在真实配对的 EEG 数据上执行。Shuffled EEG（打乱脑电）和 random EEG 只用于测试阶段的对照评估，不参与训练。这一训练与评估相分离的策略是实验设计的核心：若 shuffled/random EEG 也被纳入训练，模型可能学会利用 EEG 中的任意噪声而非真正的语义信息。仅用 real EEG 训练、同时用全部四类 EEG 模式进行评估，可以确保模型所习得的 EEG 利用能力完全源自真实配对的神经响应。

训练采用 AdamW 优化器，初始学习率设为 $1\times10^{-4}$，权重衰减为 0.01。批大小为 64（在单 GPU 上通过梯度累积等效实现），采用 cosine 学习率衰减策略，最低学习率降至初始值的 $1/100$。为缓解小样本场景下的过拟合问题，EEG encoder 中的 dropout rate 设为 0.2，gated fusion MLP 中加入了 LayerNorm 正则化。总共训练 80 个 epoch，依据验证集上的 top-1 accuracy 进行早停，patience 设为 25 个 epoch。

A2 主语义融合模型以交叉熵分类损失为主要优化目标。EEG-to-CLIP 的 InfoNCE/MSE 对齐训练主要用于前期的对比实验或可选的预训练阶段，并非最终 A2 语义融合模型的唯一训练目标。标准的训练流程分为两个阶段。第一阶段是可选的 EEG-to-CLIP 预对齐：在冻结的 CLIP 语义空间中，借助 InfoNCE 对比损失将 EEG embedding 推向对应图像 CLIP embedding 附近，为后续融合训练提供较好的初始化权重。第二阶段是语义融合训练（核心阶段）：加载对齐阶段的 checkpoint（或从随机初始化开始），联合训练 gated fusion 模块，并使用交叉熵分类损失进行优化。在第二阶段中，CLIP ViT 和 CLIP text encoder 始终保持完全冻结，仅更新 A2 EEG encoder、gated fusion 及温度参数 $\tau$。本文第 8 章所报告的 A2 结果，均来自以交叉熵损失为主训练目标、辅以可选 InfoNCE 预对齐的语义融合模型。

\subsection{模型验证方法}

为了客观判断 EEG 中是否包含可被机器利用的语义信息，A2 的评估设置了多重对照条件：real EEG 使用被试观看图像时同步记录的脑电信号；shuffled EEG（打乱脑电）将同一被试内的 EEG trial 与图像随机重新配对，破坏时间同步但保留 EEG 的统计特性；random EEG 从高斯分布中随机采样，不含任何真实的神经语义成分。只有当 real EEG 在各项指标上系统地优于 shuffled 和 random EEG 时，才能认定 EEG 确实提供了有效的语义辅助。这一对照思路贯穿了本课程设计所有实验的评估过程。

\subsection{小结}

A2 temporal-spectral-spatial EEG encoder 的核心特性可归纳如下：三个并行编码分支分别捕捉 EEG 的时间、频谱和空间信息；gated residual fusion 依据输入质量对 EEG 的辅助强度进行动态调节；CLIP 文本原型提供了约束明确的语义分类空间，避免了小样本下从头训练的困难；严格的 shuffled/random EEG 对照设置保证了效果评估的客观性。详细的实验结果将在第 8 章中给出。
